% thqg_bv_construction_extensions.tex
% BV-BRST derivation of A-infinity operations from 3d HT field theory,
% loop corrections, and genus expansion via the modular bar construction.
%
% Dependencies: bv-construction.tex (Section BV_construction),
% fm-calculus.tex (Section FM_calculus),
% higher_genus_modular_koszul.tex (Vol I, modular bar).

\providecommand{\fg}{\mathfrak{g}}
\providecommand{\fh}{\mathfrak{h}}
\providecommand{\cA}{\mathcal{A}}
\providecommand{\cO}{\mathcal{O}}
\providecommand{\cM}{\mathcal{M}}
\providecommand{\cL}{\mathcal{L}}
\providecommand{\cF}{\mathcal{F}}
\providecommand{\cH}{\mathcal{H}}
\providecommand{\cV}{\mathcal{V}}
\providecommand{\cE}{\mathcal{E}}
\providecommand{\cZ}{\mathcal{Z}}
\providecommand{\cW}{\mathcal{W}}
\providecommand{\cN}{\mathcal{N}}
\providecommand{\cD}{\mathcal{D}}
\providecommand{\bC}{\mathbb{C}}
\providecommand{\bR}{\mathbb{R}}
\providecommand{\bZ}{\mathbb{Z}}
\providecommand{\bQ}{\mathbb{Q}}
\providecommand{\bN}{\mathbb{N}}
\providecommand{\barB}{\bar{B}}
\providecommand{\barBch}{\bar{B}^{\mathrm{ch}}}
\providecommand{\Omegach}{\Omega^{\mathrm{ch}}}
\providecommand{\BBar}{\operatorname{Bar}}
\providecommand{\Cobar}{\operatorname{\Omega}}
\providecommand{\Tr}{\operatorname{Tr}}
\providecommand{\Res}{\operatorname{Res}}
\providecommand{\Sym}{\operatorname{Sym}}
\providecommand{\HH}{\operatorname{HH}}
\providecommand{\CE}{\operatorname{CE}}
\providecommand{\Hom}{\operatorname{Hom}}
\providecommand{\End}{\operatorname{End}}
\providecommand{\id}{\operatorname{id}}
\providecommand{\Spec}{\operatorname{Spec}}
\providecommand{\Obs}{\operatorname{Obs}}
\providecommand{\FM}{\operatorname{FM}}
\providecommand{\Conf}{\operatorname{Conf}}
\providecommand{\Ran}{\operatorname{Ran}}
\providecommand{\h}{\hbar}
\providecommand{\ov}[1]{\overline{#1}}
\providecommand{\Ainf}{\mathsf{A}_{\infty}}
\providecommand{\Linf}{\mathsf{L}_{\infty}}
\providecommand{\SCchtop}{\mathsf{SC}^{\mathrm{ch,top}}}
\providecommand{\Aut}{\operatorname{Aut}}
\providecommand{\Stab}{\operatorname{Stab}}
\providecommand{\Def}{\operatorname{Def}}

%% ====================================================================
%%
%% CHAPTER: BV-BRST DERIVATION OF A-INFINITY OPERATIONS
%% AND THE MODULAR BAR CONSTRUCTION
%%
%% ====================================================================

\chapter{BV-BRST derivation of \texorpdfstring{$\Ainf$}{A-infinity}
operations and the genus expansion}
% label removed: chap:thqg-bv-ext

\begin{summary}[Scope and targets]
Section~\ref{sec:BV_construction} established that the BV observable
prefactorization algebra of a 3d HT theory on $\bR_t \times \bC_z$
carries higher multiplications $m_k$, defined as integrals over
$\FM_k(\bC) \times \Conf_k(\bR)$, with Stasheff identities from
Stokes' theorem on FM boundaries. That discussion
was genus-$0$ and tree-level. Two objectives here.
First (Sections~\ref{sec:thqg-bv-ext-tree-feynman}--\ref{sec:thqg-bv-ext-mk-derivation}):
derive the tree-level $m_k$ from the 3d HT action, propagator,
and Feynman diagrams on $\FM_k(\bC) \times \Conf_k(\bR)$, showing
that the quantum master equation forces every $\Ainf$ relation
into a boundary integral. Second
(Sections~\ref{sec:thqg-bv-ext-one-loop}--\ref{sec:thqg-bv-ext-stable-graph}):
incorporate loop corrections, identify the genus-$g$ curvature
$\kappa \cdot \omega_g$, prove the two-loop three-shell identity,
and show that the stable-graph amplitude expansion
reproduces the modular bar construction of Vol~I.
\end{summary}


%% ====================================================================
%% PART 1: TREE-LEVEL BV-BRST DERIVATION OF A-INFINITY OPERATIONS
%% ====================================================================

\section{Tree-level Feynman rules on
\texorpdfstring{$\FM_k(\bC) \times \Conf_k(\bR)$}{FM\_k(C) x Conf\_k(R)}}
% label removed: sec:thqg-bv-ext-tree-feynman

\begin{convention}[Sign convention for Feynman diagrams]
% label removed: conv:thqg-bv-ext-signs
Each Feynman diagram $\Gamma$ with $L$ loops and $I$ internal
fermionic lines contributes a factor $(-1)^{L+I}$. For
genus-$0$ tree diagrams, $L = 0$ and $I = 0$ (all fields
are bosonic in the HT twist), so the sign is $+1$. At higher
genus, the loop sign $(-1)^L$ combines with the fermionic-line
sign $(-1)^I$ to produce the standard Feynman sign rule.
\end{convention}

Fix a 3d HT theory on $\bR_t \times \bC_z$ satisfying hypotheses
(H1)--(H4) of Section~\ref{subsec:analytic}. The field space $\cF$
decomposes as
\begin{equation}
% label removed: eq:thqg-bv-ext-field-decomp
\cF = \cF^{\mathrm{hol}} \otimes \cF^{\mathrm{top}},
\end{equation}
where $\cF^{\mathrm{hol}}$ carries holomorphic dependence on $z \in \bC$
and $\cF^{\mathrm{top}}$ carries topological (locally constant)
dependence on $t \in \bR$. The BV action has the form
\begin{equation}
% label removed: eq:thqg-bv-ext-action
S[\Phi] = \int_{\bR \times \bC}
\bigl\langle \Phi,\, \bar\partial_z \Phi \bigr\rangle \, dt
\;+\; \sum_{n \ge 2} \frac{1}{n!}
\int_{\bR \times \bC} V_n(\Phi,\ldots,\Phi)\, dt,
\end{equation}
with kinetic operator $\bar\partial_z$ along $\bC$ (no kinetic term
in $t$, reflecting topological invariance along $\bR$), and
interaction vertices $V_n$ of degree $2-n$ in the BV grading.

\subsection{The propagator}
% label removed: subsec:thqg-bv-ext-propagator

\begin{definition}[HT propagator; \ClaimStatusProvedHere]
% label removed: def:thqg-bv-ext-propagator
The gauge-fixed propagator is
\begin{equation}
% label removed: eq:thqg-bv-ext-prop
G(z_1, z_2;\, t_1, t_2)
\;=\;
\frac{\Theta(t_1 - t_2)}{2\pi(z_1 - z_2)},
\end{equation}
where $\Theta$ is the Heaviside step function. This is the
\emph{retarded} Green function for $\bar\partial$ on
$\bC \times \bR_+$: the holomorphic factor
$1/(2\pi(z_1 - z_2))$ is the $\bar\partial$-Green kernel on
$\bC$, and the topological factor $\Theta(t_1 - t_2)$ enforces
causal (time-)ordering. In the distributional sense,
\[
\bar\partial_{z_1} G(z_1, z_2;\, t_1, t_2)
= \delta^{(2)}(z_1 - z_2)\,\Theta(t_1 - t_2),
\]
so $G$ inverts $\bar\partial_{z_1}$ on the causal half-space
$\{t_1 > t_2\}$, not on all of $\bC \times \bR$. The retarded
boundary condition (vanishing for $t_1 < t_2$) is the HT analogue
of the Feynman $i\epsilon$ prescription.
\end{definition}

\begin{remark}[Factored form]
% label removed: rem:thqg-bv-ext-factored-propagator
The propagator factors as $G = G^{\mathrm{hol}} \cdot G^{\mathrm{top}}$
with $G^{\mathrm{hol}}(z_1,z_2) = 1/(2\pi(z_1 - z_2))$ and
$G^{\mathrm{top}}(t_1,t_2) = \Theta(t_1 - t_2)$. The holomorphic
factor produces a logarithmic $1$-form
\[
\eta_{ij} := d\log(z_i - z_j)
= \frac{dz_i - dz_j}{z_i - z_j}
\;\in\; \Omega^1_{\log}(\FM_k(\bC))
\]
on the FM compactification. The topological factor produces a
cellular cochain on $\Conf_k(\bR)$: the step function $\Theta(t_i - t_j)$
is the indicator of the half-space $\{t_i > t_j\}$ in the
$E_1$-operad. Each propagator edge thus contributes a form--chain
pairing on $\FM_k(\bC) \times \Conf_k(\bR)$.
\end{remark}

\subsection{Feynman diagrams and tree amplitudes}
% label removed: subsec:thqg-bv-ext-tree-amplitudes

\begin{definition}[Feynman diagrams for $m_k$]
% label removed: def:thqg-bv-ext-feynman-diag
A \emph{tree-level Feynman diagram} $\Gamma$ contributing to the
$k$-ary operation $m_k$ is a rooted planar tree with:
\begin{enumerate}[label=(\roman*)]
\item $k$ leaves (external legs), labeled $1, \ldots, k$ in planar
 order, corresponding to operator insertions
 $a_1, \ldots, a_k \in A$;
\item internal vertices of valence $\ge 3$, each carrying a vertex
 factor $V_n$ of the appropriate arity;
\item internal edges, each carrying a propagator
 $G(z_i, z_j;\, t_i, t_j)$;
\item a single root edge (output), corresponding to the result
 of $m_k(a_1, \ldots, a_k)$.
\end{enumerate}
The set of such diagrams is denoted $\mathrm{Tree}(k)$.
\end{definition}

\begin{definition}[Tree-level amplitude]
% label removed: def:thqg-bv-ext-tree-amplitude
For $\Gamma \in \mathrm{Tree}(k)$, the \emph{amplitude} is the
differential form on $\FM_k(\bC) \times \Conf_k(\bR)$ given by
\begin{equation}
% label removed: eq:thqg-bv-ext-tree-amp
\omega_\Gamma(a_1, \ldots, a_k)
\;:=\;
\biggl(\prod_{e = (i,j) \in E(\Gamma)}
\eta_{ij}\biggr)
\;\cdot\;
\biggl(\prod_{v \in V(\Gamma)} V_{|v|}(a_{L(v)})\biggr)
\;\cdot\;
\biggl(\prod_{e = (i,j) \in E(\Gamma)}
\Theta(t_i - t_j)\biggr),
\end{equation}
where $E(\Gamma)$ is the edge set, $V(\Gamma)$ the internal vertex
set, $|v|$ the valence of vertex $v$, and $a_{L(v)}$ denotes the
tuple of insertions at the leaves descending from $v$. The
holomorphic factor is a product of logarithmic $1$-forms $\eta_{ij}$
on $\FM_k(\bC)$; the topological factor is a product of step
functions restricting to a chamber of $\Conf_k(\bR)$.
\end{definition}

\begin{proposition}[Degree count; \ClaimStatusProvedHere]
% label removed: prop:thqg-bv-ext-degree-count
The amplitude $\omega_\Gamma$ for $\Gamma \in \mathrm{Tree}(k)$
is a form of total degree $k - 2$ on $\FM_k(\bC)$ (holomorphic
contribution) plus degree $0$ on $\Conf_k(\bR)$ (topological
contribution: the step-function product is a $0$-cochain).
A tree with $k$ leaves has $k - 1$ internal edges and hence
$k - 1$ logarithmic $1$-forms $\eta_{ij}$; the form degree is
$k - 1$ minus the one form ``used'' for the output differential,
yielding $k - 2$ on the quotient $\FM_k(\bC)/\bC$ of dimension
$2(k-1) - 2 = 2k - 4$ (real).

The operation $m_k$ has cohomological degree $2 - k$:
one unit from each internal edge ($-1$ per desuspension) plus
$+1$ from each internal vertex, yielding
$-(k-1) + (\text{number of internal vertices})$. For a
binary tree, the vertex count is $k - 1$, giving degree
$-(k-1) + (k-1) = 0$; the actual degree $2-k$ arises from
the grading shift in the bar desuspension.
\end{proposition}

\begin{proof}
Count edges and vertices. A planar rooted tree with $k$ leaves
has $k - 1$ internal edges (every non-root leaf contributes one
edge on the path to the root, and no edge is counted twice for a
tree). Each edge contributes a $1$-form $\eta_{ij}$ of
holomorphic degree $1$. The total form degree on $\FM_k(\bC)$
is therefore $k - 1$. After dividing by the translation action on
$\bC$ (which removes one complex coordinate, hence two real
dimensions), the effective degree is $k - 1$ on a space of real
dimension $2(k - 1) - 2$.

For the bar-complex degree: in the unsuspended complex, $m_k$
maps $k$ inputs to $1$ output preserving total degree, so
$|m_k|_{\mathrm{unsuspended}} = 0$. In the bar complex,
each input is desuspended by $s^{-1}$ (degree shift $-1$),
and the coderivation $d_B$ extending $m_k$ has degree $+1$
(the bar differential has cohomological degree~$+1$).
The operation $m_k$ on the bar complex has degree
$|m_k|_{\mathrm{bar}} = 2 - k$: this is because $m_k$ maps
$(s^{-1}A)^{\otimes k} \to s^{-1}A$, and
$|s^{-1}| = -1$ contributes $-k$ on the input and $-1$ on the
output, giving a net shift of $-k - (-1) = -(k-1)$; combined
with the coderivation lift (which adds $+1$), the total degree
of the coderivation component extending $m_k$ is $2 - k$.
\end{proof}

\subsection{The quantum master equation and closed weight forms}
% label removed: subsec:thqg-bv-ext-qme

\begin{proposition}[QME implies closedness; \ClaimStatusProvedHere]
% label removed: prop:thqg-bv-ext-qme-closedness
The quantum master equation
\begin{equation}
% label removed: eq:thqg-bv-ext-qme
\h\,\Delta_{\mathrm{BV}} S
+ \tfrac{1}{2}\{S, S\}_{\mathrm{BV}} = 0
\end{equation}
implies that the total differential $D = d_{\FM} + Q_{\mathrm{BRST}}$
annihilates the tree-level weight forms:
\[
D\,\omega_\Gamma = 0
\quad\text{on}\quad
\Conf_k(\bC) \times \Conf_k(\bR),
\]
for every $\Gamma \in \mathrm{Tree}(k)$.
\end{proposition}

\begin{proof}
The total differential decomposes as $D = d_{\FM} + Q$ where
$d_{\FM}$ acts on the geometric coordinates of $\FM_k(\bC)$ and
$Q = Q_{\mathrm{BRST}} = \{S, -\}_{\mathrm{BV}} + \h\,\Delta_{\mathrm{BV}}$
acts on the field insertions. At tree level ($\h = 0$), the QME
reduces to $\{S, S\}_{\mathrm{BV}} = 0$, the classical master
equation.

\medskip\noindent\textbf{Step 1: $d_{\FM}$ produces propagator
residues.}
Each logarithmic $1$-form $\eta_{ij} = d\log(z_i - z_j)$ is
closed: $d_{\FM}\,\eta_{ij} = 0$. Therefore
$d_{\FM}\,\omega_\Gamma = 0$ in the interior $\Conf_k(\bC)$, since
$\omega_\Gamma$ is a product of closed forms.

\medskip\noindent\textbf{Step 2: $Q$ produces vertex relations.}
The BRST operator $Q$ acts on the insertions $a_1, \ldots, a_k$
via the Leibniz rule. On the vertex factor $V_n(a_{i_1}, \ldots,
a_{i_n})$, the relation $Q \circ V_n = V_n \circ Q^{\otimes n}$
(a consequence of $\{S, S\} = 0$: the interaction vertices are
$Q$-equivariant) ensures that
\[
Q\,\omega_\Gamma(a_1, \ldots, a_k)
= \sum_{j=1}^{k} \pm\,
\omega_\Gamma(a_1, \ldots, Q\,a_j, \ldots, a_k).
\]
This is the ``$Q$ passes through vertices'' identity.

\medskip\noindent\textbf{Step 3: Total closure.}
Combining Steps 1--2, $D\,\omega_\Gamma = d_{\FM}\,\omega_\Gamma
+ Q\,\omega_\Gamma = 0 + Q\,\omega_\Gamma$. The total weight form
$\omega_k = \sum_\Gamma |\Aut(\Gamma)|^{-1}\,\omega_\Gamma$ is
$D$-closed if and only if the sum of all $Q$-variations vanishes,
which is guaranteed by the classical master equation applied at each
vertex. This is the content of $\{S, S\} = 0$ expanded in Feynman
graphs.
\end{proof}

\begin{remark}[Role of the QME factor $\tfrac{1}{2}$]
% label removed: rem:thqg-bv-ext-qme-factor
The factor $\tfrac{1}{2}$ in $\tfrac{1}{2}\{S,S\} = 0$ matches the
automorphism factor $|\Aut(\Gamma)|^{-1} = \tfrac{1}{2}$ for the
unique graph with two vertices joined by a single edge. More
generally, the $n!$ in the interaction $\frac{1}{n!}V_n$ compensates
the automorphism counts of tree diagrams, ensuring that the total
tree-level amplitude $\omega_k = \sum_\Gamma |\Aut(\Gamma)|^{-1}
\omega_\Gamma$ is correctly normalized.
\end{remark}


%% ====================================================================
\section{Derivation of \texorpdfstring{$m_k$}{m\_k} as a sum over
trees}
% label removed: sec:thqg-bv-ext-mk-derivation

\subsection{The tree-level operations}
% label removed: subsec:thqg-bv-ext-mk-tree

\begin{construction}[$m_k$ from Feynman diagrams;
\ClaimStatusProvedHere]
% label removed: const:thqg-bv-ext-mk-feynman
Define the $k$-ary operation $m_k : A^{\otimes k} \to A$ by
\begin{equation}
% label removed: eq:thqg-bv-ext-mk-def
m_k(a_1, \ldots, a_k)
\;:=\;
\sum_{\Gamma \in \mathrm{Tree}(k)}
\frac{1}{|\Aut(\Gamma)|}
\int_{\Gamma_k \subset \FM_k(\bC) \times \Conf_k(\bR)}
\omega_\Gamma(a_1, \ldots, a_k),
\end{equation}
where $\Gamma_k$ is the fundamental chain (the closure of the
ordered chamber $\{t_1 > \cdots > t_k\}$ in
$\FM_k(\bR) \times \FM_k(\bC)$) and
$\omega_\Gamma$ is the amplitude of
Definition~\ref{def:thqg-bv-ext-tree-amplitude}.
\end{construction}

\begin{theorem}[Tree-level $\Ainf$ from BV; \ClaimStatusProvedHere]
% label removed: thm:thqg-bv-ext-tree-ainfty
The operations $\{m_k\}_{k \ge 1}$ of
Construction~\ref{const:thqg-bv-ext-mk-feynman} satisfy the
$\Ainf$ Stasheff identities:
\begin{equation}
% label removed: eq:thqg-bv-ext-stasheff
\sum_{\substack{i + j = n + 1 \\ i, j \ge 1}}
\sum_{s=0}^{n-j}
(-1)^{|a_1| + \cdots + |a_s|}
\,m_i\bigl(a_1, \ldots, a_s,\,
m_j(a_{s+1}, \ldots, a_{s+j}),\,
a_{s+j+1}, \ldots, a_n\bigr)
= 0.
\end{equation}
\end{theorem}

\begin{proof}
Apply Stokes' theorem on $\FM_n(\bC)$ to the total weight form
$\omega_n = \sum_\Gamma |\Aut(\Gamma)|^{-1}\,\omega_\Gamma$.
By Proposition~\ref{prop:thqg-bv-ext-qme-closedness},
$D\,\omega_n = 0$ in the interior.
Theorem~\ref{thm:Stokes_FM} gives
\[
0 = \int_{\Gamma_n} D\,\omega_n
= \sum_{\substack{S \subseteq \{1,\ldots,n\} \\ |S| \ge 2}}
\epsilon_{D_S}
\int_{\Gamma_n \cap D_S} \Res_{D_S}(\omega_n).
\]
We identify each boundary integral with a composed operation.

\medskip\noindent\textbf{Step 1: Boundary factorization of
tree amplitudes.}
On the divisor $D_S \cong \FM_r(\bC) \times \FM_{|S|}(\bC)$
for a consecutive block $S = \{s+1, \ldots, s+\ell\}$, the
tree amplitude $\omega_\Gamma$ decomposes. The tree $\Gamma$
restricts to two sub-trees: an \emph{inner tree} $\Gamma_{\mathrm{in}}$
on the leaves in $S$ (contributing to $m_\ell$) and an \emph{outer
tree} $\Gamma_{\mathrm{out}}$ on the remaining leaves plus one new
leaf for the output of $\Gamma_{\mathrm{in}}$ (contributing to
$m_r$ with $r = n - \ell + 1$). At the propagator level, the
edge connecting $\Gamma_{\mathrm{in}}$ to $\Gamma_{\mathrm{out}}$
produces the logarithmic residue:
\[
\Res_{D_S}\bigl(\eta_{ij}\bigr)
= \Res_{\varepsilon_S = 0}
 \frac{d\varepsilon_S}{\varepsilon_S} = 1,
\]
where $\varepsilon_S$ is the radial coordinate of the cluster $S$
and $\eta_{ij} = d\log(z_i - z_j)$ for the edge crossing the
cluster boundary. The remaining propagators split cleanly into
inner and outer factors:
\begin{align}
% label removed: eq:thqg-bv-ext-tree-factorize
\Res_{D_S}(\omega_\Gamma)
&= \omega_{\Gamma_{\mathrm{out}}}
 \;\wedge\; \omega_{\Gamma_{\mathrm{in}}}.
\end{align}

\medskip\noindent\textbf{Step 2: Summation over trees.}
Summing~\eqref{eq:thqg-bv-ext-tree-factorize} over all
$\Gamma \in \mathrm{Tree}(n)$ and integrating, the inner
integral produces $m_\ell(a_{s+1}, \ldots, a_{s+\ell})$ and
the outer integral produces
$m_r(a_1, \ldots, a_s,\, \bullet,\, a_{s+\ell+1}, \ldots, a_n)$.
Every splitting of a tree in $\mathrm{Tree}(n)$ into inner and
outer sub-trees is in bijection with pairs
$(\Gamma_{\mathrm{in}}, \Gamma_{\mathrm{out}}) \in
\mathrm{Tree}(\ell) \times \mathrm{Tree}(r)$, and the
automorphism factors multiply:
$|\Aut(\Gamma)| = |\Aut(\Gamma_{\mathrm{in}})| \cdot
|\Aut(\Gamma_{\mathrm{out}})|$.
Therefore
\[
\sum_\Gamma \frac{1}{|\Aut(\Gamma)|}
\int_{\Gamma_n \cap D_S} \Res_{D_S}(\omega_\Gamma)
= m_r(a_1, \ldots, a_s,\,
 m_\ell(a_{s+1}, \ldots, a_{s+\ell}),\,
 a_{s+\ell+1}, \ldots, a_n).
\]

\medskip\noindent\textbf{Step 3: Non-consecutive subsets vanish.}
For non-consecutive $S$, the planar ordering of the tree forces
$\Res_{D_S}(\omega_\Gamma) = 0$: no edge of a planar tree
connects non-adjacent leaves, so the angular integral over the
cluster coordinate $\theta_S$ encounters no pole and vanishes
(cf.\ Remark~\ref{rem:nonconsecutive_vanishing}).

\medskip\noindent\textbf{Step 4: Signs.}
The orientation sign $\epsilon_{D_S}$ from
Proposition~\ref{prop:general_sign_formula} combines with the
Koszul sign from passing the inner operation $m_\ell$ (of degree
$2 - \ell$) past $a_1, \ldots, a_s$ to produce the Stasheff sign
$(-1)^{|a_1| + \cdots + |a_s|}$ in~\eqref{eq:thqg-bv-ext-stasheff}.

\medskip\noindent\textbf{Step 5: Corner cancellation.}
Codimension-$2$ corners from disjoint clusters
$D_S \cap D_T$ cancel by the Arnold--Orlik--Solomon relation
(Lemma~\ref{lem:Arnold_cancellation}): the two orderings of
residue extraction produce opposite signs. Nested corners
($S \subset T$) are sub-faces of $D_T$ accounted for by
recursive Stokes on $\FM_{|T|}(\bC)$.

\medskip\noindent\textbf{Assembly.}
Combining Steps 1--5, the Stokes identity reproduces the Stasheff
relation~\eqref{eq:thqg-bv-ext-stasheff} term by term.
\end{proof}


\subsection{Explicit tree-level operations at low arity}
% label removed: subsec:thqg-bv-ext-low-arity

\begin{example}[$m_1$: the BRST differential]
% label removed: ex:thqg-bv-ext-m1
At arity $k = 1$, $\mathrm{Tree}(1)$ contains a single diagram
with no internal edges: the identity graph. Therefore
$m_1 = Q_{\mathrm{BRST}}$, the linearized BRST operator of the
theory. This is the BV differential acting on the observable
algebra $A = \Obs(D \times I)$.
\end{example}

\begin{example}[$m_2$: the OPE]
% label removed: ex:thqg-bv-ext-m2
At arity $k = 2$, $\mathrm{Tree}(2)$ contains a single binary
vertex. The amplitude is
\[
\omega_2(a_1, a_2)
= \eta_{12} \cdot V_2(a_1, a_2) \cdot \Theta(t_1 - t_2)
= \frac{dz_1 - dz_2}{z_1 - z_2}
\cdot V_2(a_1, a_2) \cdot \Theta(t_1 - t_2).
\]
Integrating over $\FM_2(\bC) \times \{t_1 > t_2\}$:
\[
m_2(a_1, a_2)
= \int_{\FM_2(\bC)}
\frac{dz_1 - dz_2}{z_1 - z_2}
\cdot V_2(a_1, a_2)
= V_2(a_1, a_2) \cdot
 \Res_{z_1 = z_2} \frac{1}{z_1 - z_2}.
\]
This is the operator product expansion: $m_2(a_1, a_2)$ extracts
the residue of the binary OPE of $a_1$ and $a_2$.
\end{example}

\begin{example}[$m_3$: two diagrams]
% label removed: ex:thqg-bv-ext-m3
At arity $k = 3$, $\mathrm{Tree}(3)$ contains:
\begin{enumerate}[label=(\alph*)]
\item the \emph{cubic vertex}: a single trivalent vertex with
 three leaves, amplitude
 $\omega_{\mathrm{cub}} = \eta_{12} \wedge \eta_{23}
 \cdot V_3(a_1, a_2, a_3) \cdot \Theta(t_1 > t_2 > t_3)$;
\item the \emph{binary tree} (two possible shapes,
 left-branching and right-branching), each with two binary
 vertices and one internal edge. The left-branching tree has
 amplitude
 $\omega_{\mathrm{left}} = \eta_{12} \cdot \eta_{u3}
 \cdot V_2(a_1, a_2) \cdot V_2(\bullet, a_3)
 \cdot \Theta(t_1 > t_2)\,\Theta(t_u > t_3)$,
 where $u$ is the internal vertex receiving the output of $V_2(a_1, a_2)$.
\end{enumerate}
The cubic vertex contributes to $m_3$ directly; the binary trees
contribute the ``homotopy transfer'' correction ensuring that
$m_3$ satisfies the arity-$3$ Stasheff identity. On the FM boundary
$D_{\{1,2\}} \cong \FM_2(\bC) \times \FM_2(\bC)$, the cubic
vertex restricts to the OPE composition $m_2(m_2(a_1, a_2), a_3)$
and the binary tree restricts to a regular (non-singular) form.
The Stokes identity then reads: the $D_{\{1,2\}}$ and $D_{\{2,3\}}$
boundary contributions (from OPE compositions) plus the
$D_{\{1,2,3\}}$ contribution ($m_1 \circ m_3$ and
$m_3 \circ m_1^{\otimes 3}$ terms) sum to zero.
\end{example}

\subsection{The binary tree expansion and operadic composition}
% label removed: subsec:thqg-bv-ext-binary-expansion

\begin{proposition}[Binary tree expansion; \ClaimStatusProvedHere]
% label removed: prop:thqg-bv-ext-binary-expansion
For a theory with only cubic vertices ($V_n = 0$ for $n \ge 4$),
the operation $m_k$ receives contributions only from binary planar
trees. The number of such trees is the Catalan number
$C_{k-1} = \frac{1}{k}\binom{2(k-1)}{k-1}$. Each binary tree
$\Gamma$ with $k$ leaves has $k - 1$ internal edges and $k - 2$
internal vertices, and the amplitude $\omega_\Gamma$ is a product
of $k - 1$ logarithmic $1$-forms $\eta_{ij}$.
\end{proposition}

\begin{proof}
A binary tree with $k$ leaves has $k - 1$ internal edges (standard
combinatorics: each leaf except the root contributes one edge) and
$k - 2$ trivalent vertices (the Euler relation $V - E + 1 = 0$ for
a tree gives $V = E - 1 = k - 2$). Each internal edge carries a
propagator $\eta_{ij}$; there are no higher-valence vertices to
contribute additional forms. The count of binary planar trees is
the Catalan number by the ballot-problem bijection.
\end{proof}

\begin{remark}[General interaction vertices]
% label removed: rem:thqg-bv-ext-general-vertices
When higher vertices $V_n$ ($n \ge 4$) are present (as in
Virasoro, where $V_4 \ne 0$ from the nonlinear OPE
$T(z)T(w) \sim \frac{c/2}{(z-w)^4} + \cdots$, or
$\cW$-algebras), the set $\mathrm{Tree}(k)$ includes trees with
vertices of valence $\ge 4$. The Stasheff identity still holds:
it is a consequence of Stokes on $\FM_k(\bC)$ regardless of
vertex valence, since $D\,\omega_\Gamma = 0$ follows from the
QME for all tree topologies. The complication is purely
combinatorial: the tree-level amplitude is a sum over all
planar rooted trees, not just binary ones.
\end{remark}

\subsection{From tree amplitudes to the bar complex}
% label removed: subsec:thqg-bv-ext-trees-to-bar

\begin{theorem}[Tree amplitudes $=$ bar differential;
\ClaimStatusProvedHere]
% label removed: thm:thqg-bv-ext-trees-bar
Let $B(A) = \bigoplus_{k \ge 0} (sA)^{\otimes k}$ denote the
bar complex of $A$ with the bar differential
$b = \sum_{k} \sum_{s} (-1)^{\epsilon(s,j)}\,
\mathbf{1}^{\otimes s} \otimes m_j \otimes
\mathbf{1}^{\otimes(n-s-j)}$.
Then the tree-level Feynman expansion gives precisely $b^2 = 0$:
\begin{enumerate}[label=(\roman*)]
\item Each bar-complex generator $(sa_1) \otimes \cdots \otimes
 (sa_k)$ corresponds to a configuration of $k$ operator
 insertions on $\bC$ with times $t_1 > \cdots > t_k$ on $\bR$;
\item The bar differential $b$ acts by summing over all consecutive
 sub-blocks $\{s+1, \ldots, s+\ell\}$ and applying
 $m_\ell$ to that block, corresponding to the boundary of $\FM_k(\bC)$
 at the divisor $D_S$;
\item $b^2 = 0$ is Stokes' theorem:
 $\partial^2\FM_k(\bC) = 0$.
\end{enumerate}
\end{theorem}

\begin{proof}
This is the content of Theorem~\ref{thm:stokes_arnold} recast in
bar-complex language. The chain $\Gamma_k \subset
\FM_k(\bC) \times \FM_k(\bR)$ is the bar-complex generator.
The boundary $\partial\Gamma_k$ decomposes into faces
$\Gamma_r \times \Gamma_\ell$ (for $r + \ell = k + 1$) indexed
by consecutive sub-blocks, which are exactly the terms of $b$.
The identity $\partial^2 = 0$ on chains translates to $b^2 = 0$
via the residue-to-composition identification of
Section~\ref{sec:FM_calculus}.
\end{proof}

\begin{corollary}[Bar differential $=$ C-direction
factorization; \ClaimStatusProvedHere]
% label removed: cor:thqg-bv-ext-bar-c-direction
The bar differential $b$ encodes factorization in the holomorphic
($\bC$) direction:
\[
b\bigl[(sa_1) \otimes \cdots \otimes (sa_k)\bigr]
= \sum_{S \text{ consec.}} \epsilon_{D_S}\,
\Res_{D_S}\bigl[\text{FM integral of }
(sa_1) \otimes \cdots \otimes (sa_k)\bigr].
\]
The bar coproduct $\Delta$ (the coassociative coproduct on $B(A)$)
similarly encodes factorization in the topological ($\bR$)
direction: splitting the ordered sequence
$(sa_1, \ldots, sa_k)$ into two subsequences corresponds to
separating the $\bR$-coordinates into two disjoint intervals.
Together, $b$ and $\Delta$ constitute a Swiss-cheese algebra
structure on $\FM_k(\bC) \times \Conf_k(\bR)$.
\end{corollary}

\begin{proof}
The identification of $b$ with $\bC$-factorization is
Theorem~\ref{thm:thqg-bv-ext-trees-bar}(ii).
The identification of $\Delta$ with $\bR$-factorization is the
deconcatenation coproduct on the tensor coalgebra: for
$(sa_1) \otimes \cdots \otimes (sa_k)$,
\[
\Delta\bigl[(sa_1) \otimes \cdots \otimes (sa_k)\bigr]
= \sum_{p=0}^{k}
\bigl[(sa_1) \otimes \cdots \otimes (sa_p)\bigr]
\otimes
\bigl[(sa_{p+1}) \otimes \cdots \otimes (sa_k)\bigr],
\]
corresponding to the partition of $\bR$ into
$(-\infty, t_p) \cup (t_p, +\infty)$. The compatibility
$b \circ \Delta = \Delta \circ b$ (the coderivation property)
follows from the factored form of the total differential
$D = d_{\FM_k(\bC)} + d_{\Conf_k(\bR)} + Q$ on the product
space.
\end{proof}


%% ====================================================================
%% PART 2: LOOP CORRECTIONS AND GENUS EXPANSION
%% ====================================================================

\section{One-loop corrections and the curvature
\texorpdfstring{$\kappa$}{kappa}}
% label removed: sec:thqg-bv-ext-one-loop

At genus $g \geq 1$, the tree-level differential $d_0$ acquires
loop corrections from self-edge contractions, producing the full
differential $D = d_0 + \sum_{g \geq 1} \hbar^g\, d^{(g)}$.
The passage from tree-level to loop-level is the BV-theoretic
incarnation of the genus expansion in the modular bar construction.

\medskip
Tree-level amplitudes compute the genus-$0$ $\Ainf$ structure.
At one loop, the BV Laplacian $\Delta_{\mathrm{BV}}$ contributes, and
the quantum master equation~\eqref{eq:thqg-bv-ext-qme} forces a
curvature correction.

\subsection{One-loop Feynman diagrams}
% label removed: subsec:thqg-bv-ext-one-loop-diagrams

\begin{definition}[One-loop diagrams]
% label removed: def:thqg-bv-ext-one-loop-diag
A \emph{one-loop Feynman diagram} $\Gamma$ contributing to the
genus-$1$ correction has:
\begin{enumerate}[label=(\roman*)]
\item $k$ external legs ($k \ge 0$);
\item exactly one cycle (loop number $g = 1$), so
 $|E(\Gamma)| - |V(\Gamma)| + 1 = 1$;
\item internal edges carrying propagators
 $G(z_i, z_j;\, t_i, t_j)$ and vertices carrying $V_n$.
\end{enumerate}
The simplest one-loop diagram has $k = 0$ external legs and a
single edge forming a self-loop at one vertex: this is the
\emph{tadpole}, whose amplitude involves the coincident-point
limit $G(z, z;\, t, t)$.
\end{definition}

\begin{proposition}[One-loop amplitude; \ClaimStatusProvedHere]
% label removed: prop:thqg-bv-ext-one-loop-amp
The one-loop diagram with $k$ external legs and one internal loop
has amplitude
\begin{equation}
% label removed: eq:thqg-bv-ext-one-loop-amp
\omega_\Gamma^{(1)}(a_1, \ldots, a_k)
= \biggl(\prod_{e \in E_{\mathrm{tree}}} \eta_e \biggr)
\cdot \eta_{\mathrm{loop}}
\cdot \biggl(\prod_v V_{|v|}(a_{L(v)})\biggr)
\cdot \biggl(\prod_e \Theta_e\biggr),
\end{equation}
where $E_{\mathrm{tree}}$ denotes the tree edges (forming a
spanning tree of $\Gamma$), $\eta_{\mathrm{loop}}$ is the
logarithmic $1$-form on the loop edge, and the product of
$\Theta$-functions enforces time-ordering along the loop as well
as on the tree edges.

The loop integral requires regularization. On a genus-$1$ surface
$\Sigma_1$ (obtained by cutting the loop and sewing to a torus),
the propagator $G$ is replaced by the genus-$1$ prime form. The
regulated amplitude converges and produces a curvature term.
\end{proposition}

\begin{proof}
The edge and vertex count follows from the definition. A graph
with one loop has $|E| = |V|$, so $k$ external legs, $|V|$
internal vertices (each with valence $\ge 3$), and $|V|$ internal
edges. One edge closes the loop; the remaining $|V| - 1$ form a
spanning tree. The amplitude is the product of propagator forms
over all edges, vertex factors, and time-ordering step functions.
Regularization on $\Sigma_1$ replaces the flat-space propagator
$1/(z_1 - z_2)$ by the Szego kernel $S_1(z_1, z_2|\tau)$, which
has the same OPE singularity but is doubly periodic. The
convergence of the resulting integral is guaranteed by
hypothesis (H1c) (one-loop finiteness).
\end{proof}

\subsection{The curvature $\kappa$ from the vacuum one-loop diagram}
% label removed: subsec:thqg-bv-ext-kappa

\begin{theorem}[Genus-$1$ curvature; \ClaimStatusProvedHere]
% label removed: thm:thqg-bv-ext-kappa
The vacuum one-loop amplitude ($k = 0$) produces the curvature
scalar
\begin{equation}
% label removed: eq:thqg-bv-ext-kappa-def
\kappa(A)
= \Tr_{A}\bigl((-1)^F q^{L_0}\bigr)\big|_{q^1}
= \sum_{n} (-1)^{|a_n|} \langle a_n,\, G \cdot a_n \rangle,
\end{equation}
where the trace is over a basis $\{a_n\}$ of $A$ and $G$ is the
propagator contracted along the loop. In terms of the conformal
weight $h$ and BV grading, $\kappa$ equals the supertrace of
$L_0$ on $A$:
\begin{equation}
% label removed: eq:thqg-bv-ext-kappa-supertrace
\kappa(A) = \operatorname{str}_{A}(L_0).
\end{equation}

The genus-$1$ fibrewise bar differential acquires a curvature:
\begin{equation}
% label removed: eq:thqg-bv-ext-curved-bar
b_1^2 = \kappa \cdot \omega_1,
\end{equation}
where $\omega_1 \in H^1(\overline{\cM}_{1,1})$ is the fundamental
class of the moduli space of elliptic curves and $b_1$ denotes the
bar differential extended to genus $1$ by including the one-loop
amplitude.
\end{theorem}

\begin{proof}
\textbf{Step 1: The tadpole.}
The vacuum one-loop diagram is a single vertex with a self-loop.
Its amplitude is
\[
\omega_{\mathrm{tad}}
= \oint_{\gamma} \frac{dz}{2\pi(z - w)}\bigg|_{w \to z}
\cdot V_2(a_n, a^n),
\]
where $\{a_n\}$ is a basis and $\{a^n\}$ its dual. On a flat
torus $\bC / (\bZ + \bZ\tau)$, the propagator is replaced by the
Weierstrass $\zeta$-function (the logarithmic derivative of the
$\sigma$-function), and the integral picks up the quasi-period
$\eta_1(\tau)$. Summing over the basis of $A$ with BV signs
produces $\kappa = \operatorname{str}(L_0)$.

\medskip\noindent\textbf{Step 2: Curvature of the bar complex.}
At genus $1$, the bar differential includes both tree-level
compositions (from $\partial\FM_k(\bC)$) and the one-loop
correction (from the BV Laplacian $\Delta_{\mathrm{BV}}$).
The QME $\h\,\Delta_{\mathrm{BV}}S + \tfrac{1}{2}\{S,S\} = 0$
now contributes at order $\h^1$:
\[
D^2 = \h\,\Delta_{\mathrm{BV}}\bigl(\{S, -\}\bigr)
= \h\,\kappa \cdot \omega_1 \ne 0.
\]
This is the genus-$1$ curvature of the curved $\Ainf$ algebra:
$m_1^2(a) = [\kappa \cdot \omega_1,\, a]$ (commutator, with the
minus sign from the BV bracket). The curvature $\kappa \cdot
\omega_1$ is a scalar multiple of the Mumford class on
$\overline{\cM}_{1,1}$.

\medskip\noindent\textbf{Step 3: Identification with
Vol~I formula.}
The supertrace formula $\kappa = \operatorname{str}(L_0)$ matches
the Vol~I definition (see the modular characteristic of
Theorem~D): $\kappa(A) = \sum_n (-1)^n \cdot n \cdot
\dim A_n$, where $A_n$ is the weight-$n$ subspace. For the
Heisenberg algebra $\cH_k$, $\kappa = k$; for affine
$\hat{\fg}_k$, $\kappa = \dim\fg \cdot (k + h^\vee) / (2h^\vee)$; for the
Virasoro algebra, $\kappa = c/2$.
\end{proof}

\begin{remark}[Higher-loop contributions to the curvature]
% label removed: rem:thqg-bv-ext-kappa-higher-loop
The curvature $\kappa$ is determined by the genus-$1$ tadpole
(one-loop vacuum diagram) alone. Higher-loop vacuum diagrams
on $\overline{\cM}_{g,0}$ evaluate to $\kappa^g \cdot \lambda_g$
by the genus universality theorem (Theorem~D, Vol~I): there are
no \emph{additional} contributions to the curvature beyond the
one-loop tadpole. The higher-genus amplitudes $\cA^{(g)}_0$
are determined recursively from $\kappa$ and the tree-level data
by the MC recursion, but they do not modify the value of $\kappa$
itself. This is because $\kappa$ is a genus-$0$-to-genus-$1$
obstruction class, and the recursion at genus $g \geq 2$ solves
for $\Theta^{(g)}$ in terms of $\kappa$ and the lower-genus
elements without feeding back.
\end{remark}

\begin{remark}[Curved $\Ainf$ vs.\ flat $\Ainf$]
% label removed: rem:thqg-bv-ext-curved-flat
At genus $0$ the bar complex has $b^2 = 0$ (flat). At genus
$g \ge 1$, the one-loop curvature $\kappa \cdot \omega_g$ makes
the bar complex \emph{curved}: $b_g^2 = \kappa \cdot \omega_g
\ne 0$ generically. The passage from flat to curved is the
BV-theoretic avatar of the passage from the genus-$0$ to the
all-genera modular operad. The bar complex at genus $g$ lives
not in the ordinary derived category but in the
\emph{coderived} category, as enforced by the nonvanishing
curvature (cf.\ Vol~I, Section on coderived categories for
curved algebras).
\end{remark}


%% ====================================================================
\section{Two-loop corrections and the three-shell identity}
% label removed: sec:thqg-bv-ext-two-loop

\subsection{Two-loop Feynman diagrams}
% label removed: subsec:thqg-bv-ext-two-loop-diagrams

\begin{definition}[Two-loop diagrams]
% label removed: def:thqg-bv-ext-two-loop-diag
A \emph{two-loop Feynman diagram} $\Gamma$ has loop number
$g = 2$: $|E(\Gamma)| - |V(\Gamma)| + 1 = 2$. The three
simplest topologies at genus $2$ are:
\begin{enumerate}[label=(\alph*)]
\item the \emph{theta graph} $\Theta$: two vertices connected by
 three edges ($|V| = 2$, $|E| = 3$, $g = 2$);
\item the \emph{figure-eight}: one vertex with two self-loops
 ($|V| = 1$, $|E| = 2$, $g = 2$);
\item the \emph{sunset}: three vertices on a triangle with one
 external edge ($|V| = 3$, $|E| = 4$, $g = 2$).
\end{enumerate}
Each contributes to the genus-$2$ amplitude via integration over
the moduli space $\overline{\cM}_{2,k}$ of genus-$2$ curves
with $k$ marked points.
\end{definition}

\begin{proposition}[Two-loop amplitude on $\overline{\cM}_{2,k}$;
\ClaimStatusProvedHere]
% label removed: prop:thqg-bv-ext-two-loop-amp
The two-loop contribution to the $k$-point amplitude is
\begin{equation}
% label removed: eq:thqg-bv-ext-two-loop-amp
\cA_k^{(2)}(a_1, \ldots, a_k)
= \sum_{\Gamma:\, g(\Gamma) = 2}
\frac{1}{|\Aut(\Gamma)|}
\int_{\overline{\cM}_{2,k}}
\omega_\Gamma^{(2)}(a_1, \ldots, a_k),
\end{equation}
where $\omega_\Gamma^{(2)}$ is the genus-$2$ weight form built from
the genus-$2$ prime form $E_2(z_i, z_j)$ and the genus-$2$ Szego
kernel $S_2(z_i, z_j)$ replacing the flat-space propagator on each
edge.
\end{proposition}

\begin{proof}
The genus-$2$ surface $\Sigma_2$ arises from the loop structure of
$\Gamma$: a graph with two independent cycles has first Betti
number $b_1(\Gamma) = 2$, and the Feynman parametrization assigns
moduli in $\cM_2$ (period matrix, marked points). Each propagator
edge on $\Sigma_2$ carries the Szego kernel
$S_2(z_i, z_j) = \partial_{z_i}\partial_{z_j}\log E_2(z_i, z_j)$,
where $E_2$ is the genus-$2$ prime form (a section of
$K^{-1/2} \boxtimes K^{-1/2}$ on $\Sigma_2 \times \Sigma_2$).
The integration over $\overline{\cM}_{2,k}$ includes both the
vertex positions on $\Sigma_2$ and the moduli of $\Sigma_2$ itself.
Convergence follows from the compactification
$\overline{\cM}_{2,k}$ and the logarithmic behavior of the prime
form at the boundary divisors (Deligne--Mumford compactification).
\end{proof}

\subsection{The three-shell identity at genus $2$}
% label removed: subsec:thqg-bv-ext-three-shell

The boundary of $\overline{\cM}_{2,k}$ decomposes into strata
indexed by stable dual graphs of genus $2$. The codimension-$1$
boundary has three types of strata (the ``three shells''):

\begin{definition}[Three shells of $\partial\overline{\cM}_{2,k}$]
% label removed: def:thqg-bv-ext-three-shells
\begin{enumerate}[label=(\roman*)]
\item \textbf{Separating node} (shell I): $\Sigma_2$ degenerates
 into $\Sigma_{1,S} \cup \Sigma_{1,S^c}$, two genus-$1$ components
 meeting at a node. The dual graph is a single edge joining two
 genus-$1$ vertices. External legs are distributed between $S$
 and $S^c$.
\item \textbf{Non-separating node} (shell II): $\Sigma_2$
 degenerates by pinching a non-separating cycle, producing
 $\Sigma_{1,k+2}$ (a genus-$1$ surface with two extra marked
 points at the node). The dual graph is a single genus-$1$
 vertex with a self-loop.
\item \textbf{Two non-separating nodes} (shell III): $\Sigma_2$
 degenerates into $\Sigma_{0,k+4}$ (a genus-$0$ surface) by
 pinching two non-separating cycles. The dual graph is a single
 genus-$0$ vertex with two self-loops.
\end{enumerate}
\end{definition}

\begin{theorem}[Three-shell identity; \ClaimStatusProvedHere]
% label removed: thm:thqg-bv-ext-three-shell
Stokes' theorem on $\overline{\cM}_{2,k}$ applied to the genus-$2$
weight form $\omega^{(2)}_k$ yields the \emph{three-shell identity}:
\begin{equation}
% label removed: eq:thqg-bv-ext-three-shell-id
0 = \sum_{\text{shell I}} \epsilon_I\,
\int_{\partial_I \overline{\cM}_{2,k}}
\Res_I(\omega^{(2)}_k)
\;+\; \sum_{\text{shell II}} \epsilon_{II}\,
\int_{\partial_{II} \overline{\cM}_{2,k}}
\Res_{II}(\omega^{(2)}_k)
\;+\; \sum_{\text{shell III}} \epsilon_{III}\,
\int_{\partial_{III} \overline{\cM}_{2,k}}
\Res_{III}(\omega^{(2)}_k),
\end{equation}
where each residue decomposes into products of lower-genus
amplitudes:
\begin{itemize}
\item Shell I: $\Res_I = \cA_{|S|}^{(1)} \otimes
 \cA_{|S^c|}^{(1)}$ (product of two genus-$1$ amplitudes,
 contracted at the node);
\item Shell II: $\Res_{II} = \Delta_{\mathrm{sew}}\bigl(
 \cA_{k+2}^{(1)}\bigr)$ (genus-$1$ amplitude with the two
 nodal points sewn by the propagator, i.e., the sewing
 operator $\Delta_{\mathrm{sew}}$ applied to the one-loop
 amplitude);
\item Shell III: $\Res_{III} = \Delta_{\mathrm{sew}}^{(2)}\bigl(
 \cA_{k+4}^{(0)}\bigr)$ (tree-level amplitude with two pairs
 of points sewn, i.e., the double sewing operator applied to
 the genus-$0$ amplitude).
\end{itemize}
\end{theorem}

\begin{proof}
\textbf{Step 1: Stokes on $\overline{\cM}_{2,k}$.}
The moduli space $\overline{\cM}_{2,k}$ is a smooth orbifold
(Deligne--Mumford stack) with normal-crossings boundary. Its
codimension-$1$ boundary divisors are indexed by stable dual
graphs $\Gamma$ with $\sum_v g(v) + b_1(\Gamma) = 2$ and $k$
half-edges. Stokes' theorem gives
\[
0 = \int_{\overline{\cM}_{2,k}} d\,\omega^{(2)}_k
= \sum_{\Gamma \in \partial\overline{\cM}_{2,k}}
\frac{\epsilon_\Gamma}{|\Aut(\Gamma)|}
\int_{D_\Gamma} \Res_{D_\Gamma}(\omega^{(2)}_k),
\]
where $D_\Gamma \cong \prod_v \overline{\cM}_{g(v), n(v)}$ is
the stratum and $\epsilon_\Gamma$ the orientation sign.

\medskip\noindent\textbf{Step 2: Classification of strata.}
The stable dual graphs with $g = 2$ and one edge (codimension $1$)
are:
\begin{itemize}
\item Edge connecting two distinct vertices of genus $1$
 (separating node, shell I);
\item Self-loop at a vertex of genus $1$ (non-separating node,
 shell II);
\item Two self-loops at a vertex of genus $0$ (two non-separating
 nodes, shell III).
\end{itemize}
The three-shell decomposition of
Definition~\ref{def:thqg-bv-ext-three-shells} exhausts all
codimension-$1$ strata.

\medskip\noindent\textbf{Step 3: Residue identification.}
On each stratum $D_\Gamma$, the weight form restricts by the
node-factorization property of the prime form:
\[
E_2(z_i, z_j)\big|_{D_\Gamma}
= E_{g(v)}(z_i, z_j) \cdot
(\text{cross-terms at the node}).
\]
The cross-terms at the node contribute the sewing operator
$\Delta_{\mathrm{sew}} = \sum_n a_n \otimes a^n$ (the propagator
contracted across the node, summed over a basis). For shell I, the
factorization is $\cA^{(1)} \otimes \cA^{(1)}$ with a single
contraction at the node. For shell II, the self-loop produces
$\Delta_{\mathrm{sew}}$ acting on the single genus-$1$ component.
For shell III, two self-loops produce
$\Delta_{\mathrm{sew}}^{(2)}$ on the genus-$0$ component.

\medskip\noindent\textbf{Step 4: Corner cancellation.}
Codimension-$2$ corners (two simultaneous degenerations) cancel by
$\partial^2 = 0$ on $\overline{\cM}_{2,k}$, the Deligne--Mumford
analogue of the Arnold--Orlik--Solomon relation. The two
orderings of degeneration produce opposite orientation signs.

\medskip\noindent\textbf{Assembly.}
The sum of shell I, II, and III contributions equals
$\int d\omega^{(2)} = 0$, yielding~\eqref{eq:thqg-bv-ext-three-shell-id}.
\end{proof}

\begin{corollary}[Genus-$2$ curvature; \ClaimStatusProvedHere]
% label removed: cor:thqg-bv-ext-genus2-curvature
At $k = 0$ (vacuum amplitude), the three-shell identity
reduces to:
\begin{equation}
% label removed: eq:thqg-bv-ext-genus2-vac
\cA^{(2)}_{\mathrm{vac}}
= \frac{\kappa^2}{2}\,[\omega_1]^2
\;+\; \kappa\,\Delta_{\mathrm{sew}}(\omega_1)
\;+\; \Delta_{\mathrm{sew}}^{(2)}(1),
\end{equation}
where the three terms correspond to shell I ($\kappa \cdot \kappa$
from two genus-$1$ tadpoles), shell II ($\kappa$ times the sewing
of one genus-$1$ diagram), and shell III (double sewing of the
genus-$0$ identity). The quantity $\cA^{(2)}_{\mathrm{vac}}$ is
the genus-$2$ partition function of $A$, expressed entirely in
terms of genus-$0$ and genus-$1$ data plus the sewing operator.
\end{corollary}

\begin{proof}
Set $k = 0$ in the three-shell identity. Shell I with $k = 0$
has both components as vacuum genus-$1$ amplitudes, each equal to
$\kappa \cdot \omega_1$. Their product, contracted at the node,
gives $\kappa^2\,[\omega_1]^2 / 2$ (the factor $1/2$ from the
automorphism of the separating edge). Shell II with $k = 0$ has
the genus-$1$ amplitude ($\kappa$ times a weight form) with two
extra points sewn, giving $\kappa\,\Delta_{\mathrm{sew}}(\omega_1)$.
Shell III with $k = 0$ is the double sewing of the trivial genus-$0$
amplitude, giving $\Delta_{\mathrm{sew}}^{(2)}(1)$.
\end{proof}


%% ====================================================================
\section{The stable-graph expansion and the modular bar construction}
% label removed: sec:thqg-bv-ext-stable-graph

\subsection{Stable graphs and the Feynman transform}
% label removed: subsec:thqg-bv-ext-stable-graphs

\begin{definition}[Stable graphs]
% label removed: def:thqg-bv-ext-stable-graph
A \emph{stable graph} $\Gamma$ of type $(g, k)$ consists of:
\begin{enumerate}[label=(\roman*)]
\item a finite graph with vertex set $V(\Gamma)$ and edge set
 $E(\Gamma)$;
\item a genus function $g_v : V(\Gamma) \to \bZ_{\ge 0}$ assigning
 a genus to each vertex;
\item $k$ labeled half-edges (legs) attached to vertices;
\item the stability condition $2g_v - 2 + n_v > 0$ at each vertex,
 where $n_v = \mathrm{val}(v) + (\text{legs at } v)$;
\item the genus relation $g = \sum_v g_v + b_1(\Gamma)$, where
 $b_1(\Gamma) = |E(\Gamma)| - |V(\Gamma)| + 1$ is the first
 Betti number.
\end{enumerate}
The set of all stable graphs of type $(g, k)$ is denoted
$\mathrm{Graph}(g, k)$.
\end{definition}

\begin{construction}[Stable-graph amplitude]
% label removed: const:thqg-bv-ext-stable-amp
For $\Gamma \in \mathrm{Graph}(g, k)$, the \emph{amplitude} is
\begin{equation}
% label removed: eq:thqg-bv-ext-stable-amp
\cA_\Gamma(a_1, \ldots, a_k)
= \int_{\prod_v \overline{\cM}_{g_v, n_v}}
\biggl(\prod_{e \in E(\Gamma)} \Delta_e\biggr)
\cdot
\biggl(\prod_{v \in V(\Gamma)}
\omega_{g_v, n_v}(a_{L(v)})\biggr),
\end{equation}
where:
\begin{itemize}
\item $\omega_{g_v, n_v}$ is the weight form at vertex $v$,
 built from the genus-$g_v$ prime form and vertex factors;
\item $\Delta_e = \sum_n a_n \otimes a^n$ is the sewing
 operator at edge $e$, contracting the two half-edges with the
 propagator;
\item the integration is over the product of moduli spaces
 $\prod_v \overline{\cM}_{g_v, n_v}$.
\end{itemize}
The total genus-$g$, $k$-point amplitude is the sum over all
stable graphs:
\begin{equation}
% label removed: eq:thqg-bv-ext-total-amp
\cA_k^{(g)}(a_1, \ldots, a_k)
= \sum_{\Gamma \in \mathrm{Graph}(g, k)}
\frac{\h^{g}}{|\Aut(\Gamma)|}\,
\cA_\Gamma(a_1, \ldots, a_k).
\end{equation}
\end{construction}

\begin{remark}[The $\h$-expansion]
% label removed: rem:thqg-bv-ext-hbar
The loop expansion parameter $\h$ tracks the genus: a graph
$\Gamma$ of genus $g = \sum_v g_v + b_1(\Gamma)$ contributes at
order $\h^g$. The full amplitude is the formal power series
\[
\cA_k = \sum_{g \ge 0} \h^g\,\cA_k^{(g)}
\;\in\; A^{\otimes k}\llbracket\h\rrbracket.
\]
At $\h = 0$, only trees contribute ($g = 0$), recovering the
tree-level $\Ainf$ structure of
Section~\ref{sec:thqg-bv-ext-mk-derivation}.
\end{remark}

\subsection{Stokes on stable-graph strata: the general identity}
% label removed: subsec:thqg-bv-ext-stokes-stable

\begin{theorem}[Stable-graph Stokes identity; \ClaimStatusProvedHere]
% label removed: thm:thqg-bv-ext-stable-stokes
The QME combined with Stokes' theorem on
$\prod_v \overline{\cM}_{g_v, n_v}$ yields the identity
\begin{equation}
% label removed: eq:thqg-bv-ext-stable-stokes
\sum_{\Gamma' \in \partial\Gamma}
\frac{\epsilon_{\Gamma'}}{|\Aut(\Gamma')/\Aut(\Gamma)|}
\,\cA_{\Gamma'}
= 0,
\end{equation}
where $\partial\Gamma$ denotes the set of stable graphs obtained
from $\Gamma$ by:
\begin{enumerate}[label=(\alph*)]
\item \textbf{edge contraction}: contracting an edge $e$ of
 $\Gamma$, merging its two endpoints into a single vertex of
 genus $g_{v_1} + g_{v_2} + \delta_{e}$ (where
 $\delta_e = 1$ if $e$ is a self-loop, $0$ otherwise);
\item \textbf{vertex splitting}: splitting a vertex $v$ of genus
 $g_v$ into two vertices of genera $g_1, g_2$ with
 $g_1 + g_2 = g_v$ (or $g_1 + g_2 + 1 = g_v$ if a new edge is
 created), distributing the half-edges at $v$ between the two
 new vertices.
\end{enumerate}
These two operations correspond to the two types of codimension-$1$
boundary strata of $\overline{\cM}_{g,k}$: node formation
(vertex splitting) and the inverse operation (edge contraction).
\end{theorem}

\begin{proof}
Apply Stokes' theorem to the amplitude $\cA_\Gamma$ on its
integration domain $\prod_v \overline{\cM}_{g_v, n_v}$. The
boundary of each factor $\overline{\cM}_{g_v, n_v}$ consists of
divisors indexed by stable dual graphs with one additional edge
(a degeneration of the genus-$g_v$ curve). Each such degeneration
at vertex $v$ produces a new stable graph $\Gamma'$ with one more
edge than $\Gamma$, obtained by splitting $v$. The residue at the
boundary divisor factors the weight form:
\[
\Res_{D}(\omega_{g_v, n_v})
= \omega_{g_1, n_1} \cdot \Delta_{\mathrm{new}} \cdot
\omega_{g_2, n_2},
\]
where $\Delta_{\mathrm{new}} = \sum_n a_n \otimes a^n$ is the
sewing operator at the new node. Inserting this into the product
over all vertices of $\Gamma$ and summing over all possible
degenerations gives the sum over $\Gamma' \in \partial\Gamma$.

The identity $\sum \epsilon_{\Gamma'}\,\cA_{\Gamma'} = 0$ is
the chain-level statement of $\partial^2 = 0$ on the space of
stable graphs, which is the defining property of the
\emph{Feynman transform} of the modular operad. The signs
$\epsilon_{\Gamma'}$ arise from the orientation convention on
moduli spaces (the Mumford convention).

Corner cancellation (codimension-$2$ strata: two simultaneous
degenerations) follows from the normal-crossings property of
$\partial\overline{\cM}_{g,k}$ and the relation $\partial^2 = 0$.
\end{proof}

\subsection{Identification with the modular bar construction}
% label removed: subsec:thqg-bv-ext-modular-bar

\begin{theorem}[Stable-graph expansion $=$ modular bar;
\ClaimStatusProvedHere]
% label removed: thm:thqg-bv-ext-modular-bar
The genus-$g$, $k$-point amplitude $\cA_k^{(g)}$
of~\eqref{eq:thqg-bv-ext-total-amp} equals the $(g,k)$-component
of the modular bar construction $B^{(g,k)}(A)$ of Vol~I:
\begin{equation}
% label removed: eq:thqg-bv-ext-modular-bar-id
\cA_k^{(g)}(a_1, \ldots, a_k)
= B^{(g,k)}(A)(a_1, \ldots, a_k).
\end{equation}
Under this identification:
\begin{enumerate}[label=(\roman*)]
\item the stable-graph summation
 in~\eqref{eq:thqg-bv-ext-total-amp} is the graph sum defining
 the Feynman transform $\mathsf{F}(\mathsf{Com})$ of the
 commutative modular operad;
\item the sewing operator $\Delta_e$ at each edge is the
 contraction with the invariant pairing
 $\langle -,- \rangle : A \otimes A \to \bC$;
\item the weight forms $\omega_{g_v, n_v}$ at each vertex are the
 $(g_v, n_v)$-components of the chiral algebra structure
 maps of $A$;
\item the differential on $B^{(g,k)}(A)$ (the modular bar
 differential) decomposes as
 \begin{equation}
 % label removed: eq:thqg-bv-ext-modular-diff
 D_{\mathrm{mod}}
 = d_{\mathrm{int}} + d_{\mathrm{sew}} + \h\,\Delta_{\mathrm{BV}},
 \end{equation}
 where $d_{\mathrm{int}}$ acts within each vertex (the internal
 differential, from $m_1$), $d_{\mathrm{sew}}$ creates new edges
 (sewing, from boundary strata of $\overline{\cM}_{g_v, n_v}$),
 and $\h\,\Delta_{\mathrm{BV}}$ closes existing half-edges into
 self-loops (the BV Laplacian);
\item $D_{\mathrm{mod}}^2 = 0$ is the QME, which at each genus
 $g$ is the Stokes identity of
 Theorem~\ref{thm:thqg-bv-ext-stable-stokes}.
\end{enumerate}
\end{theorem}

\begin{proof}
The identification proceeds component by component.

\medskip\noindent\textbf{Genus $0$.}
At $g = 0$, the stable graphs are trees (since $b_1 = 0$ and all
vertex genera are $0$). The amplitude
$\cA_k^{(0)} = \sum_{\Gamma \in \mathrm{Tree}(k)}
|\Aut(\Gamma)|^{-1}\,\cA_\Gamma$ is the tree-level operation
$m_k$ of Construction~\ref{const:thqg-bv-ext-mk-feynman}. The
bar complex $B^{(0,k)}(A) = (sA)^{\otimes k}$ with differential
$b = \sum m_j$ is the tree-level bar complex. The identification
$\cA_k^{(0)} = B^{(0,k)}$ is
Theorem~\ref{thm:thqg-bv-ext-trees-bar}.

\medskip\noindent\textbf{Genus $1$.}
At $g = 1$, the stable graphs have either one loop edge on a
tree (contributing the BV Laplacian $\Delta_{\mathrm{BV}}$) or
a genus-$1$ vertex with no loops (contributing the genus-$1$
structure map). The differential
$D_{\mathrm{mod}} = d_{\mathrm{int}} + d_{\mathrm{sew}}
+ \h\Delta_{\mathrm{BV}}$ at genus $1$ has
$D_{\mathrm{mod}}^2 = \h\,\kappa \cdot \omega_1$: the
genus-$1$ curvature of
Theorem~\ref{thm:thqg-bv-ext-kappa}. This matches
$B^{(1,k)}(A)$ with its curved differential.

\medskip\noindent\textbf{Genus $2$.}
At $g = 2$, the three-shell identity
(Theorem~\ref{thm:thqg-bv-ext-three-shell}) decomposes
$D_{\mathrm{mod}}^2 = 0$ at genus $2$ into contributions
from the three boundary types of $\overline{\cM}_{2,k}$.
This matches $B^{(2,k)}(A)$.

\medskip\noindent\textbf{General genus.}
At arbitrary genus $g$, the argument proceeds by induction on $g$.
The stable-graph summation is the definition of
$\mathsf{F}(\mathsf{Com})$ applied to $A$. The differential
$D_{\mathrm{mod}}$ is the Feynman-transform differential, whose
square vanishes by $\partial^2 = 0$ on the modular operad
$\overline{\cM}_{g,k}$. The QME ensures compatibility between
the vertex structure maps and the sewing operators.

The automorphism factors $|\Aut(\Gamma)|^{-1}$ in the graph
sum match the orbifold multiplicities of
$\overline{\cM}_{g,k}$: each boundary stratum $D_\Gamma$ has
stabilizer $\Aut(\Gamma)$ in the Deligne--Mumford stack.
\end{proof}

\subsection{The modular bar differential as boundary operator}
% label removed: subsec:thqg-bv-ext-modular-diff

\begin{corollary}[Three components of $D_{\mathrm{mod}}$;
\ClaimStatusProvedHere]
% label removed: cor:thqg-bv-ext-three-components
The modular bar differential decomposes into three geometrically
distinct operations:
\begin{equation}
% label removed: eq:thqg-bv-ext-three-comp
D_{\mathrm{mod}} = d_{\mathrm{int}} + d_{\mathrm{sew}}
+ \h\,\Delta_{\mathrm{BV}},
\end{equation}
where:
\begin{enumerate}[label=(\roman*)]
\item $d_{\mathrm{int}} = \sum_{v} m_1^{(v)}$ acts by
 applying the BRST differential $m_1$ at each vertex
 independently; this is the \emph{internal differential},
 corresponding to the genus-$0$ BRST cohomology at each
 vertex;
\item $d_{\mathrm{sew}} = \sum_{\text{pairs}} m_2^{(v,w)}$
 creates a new edge between two half-edges at distinct
 vertices $v, w$ (or at the same vertex, increasing the
 loop count); this is the \emph{sewing operator},
 corresponding to the boundary stratum of
 $\overline{\cM}_{g_v, n_v}$ where two marked points collide;
\item $\h\,\Delta_{\mathrm{BV}} = \sum_{\text{half-edges}}
 \Delta^{(v)}$ closes a pair of half-edges at a single
 vertex into a self-loop; this is the \emph{BV Laplacian},
 corresponding to the non-separating degeneration of the
 genus-$g_v$ surface (pinching a non-separating cycle).
\end{enumerate}
The three components satisfy:
\begin{align}
% label removed: eq:thqg-bv-ext-d-squared
d_{\mathrm{int}}^2 &= 0,
&
d_{\mathrm{int}}\,d_{\mathrm{sew}}
+ d_{\mathrm{sew}}\,d_{\mathrm{int}} &= 0,\\
d_{\mathrm{sew}}^2
+ d_{\mathrm{int}}\,\h\Delta_{\mathrm{BV}}
+ \h\Delta_{\mathrm{BV}}\,d_{\mathrm{int}} &= 0,
&
d_{\mathrm{sew}}\,\h\Delta_{\mathrm{BV}}
+ \h\Delta_{\mathrm{BV}}\,d_{\mathrm{sew}} &= 0,\\
(\h\Delta_{\mathrm{BV}})^2 &= 0.
& &
\end{align}
These five relations are the $\h$-graded components of
$D_{\mathrm{mod}}^2 = 0$, expressing the QME order by order in
the genus expansion.
\end{corollary}

\begin{proof}
The decomposition~\eqref{eq:thqg-bv-ext-three-comp} follows from
Theorem~\ref{thm:thqg-bv-ext-modular-bar}(iv).
The five relations are obtained by expanding
$D_{\mathrm{mod}}^2 = (d_{\mathrm{int}} + d_{\mathrm{sew}}
+ \h\Delta_{\mathrm{BV}})^2 = 0$ and collecting terms by the
number of edges created (equivalently, by the power of $\h$):
\begin{itemize}
\item $\h^0$, no new edges: $d_{\mathrm{int}}^2 = 0$
 ($m_1^2 = 0$ at each vertex, the tree-level Stasheff identity
 at arity $1$);
\item $\h^0$, one new edge:
 $d_{\mathrm{int}}\,d_{\mathrm{sew}}
 + d_{\mathrm{sew}}\,d_{\mathrm{int}} = 0$ (compatibility
 of the BRST differential with sewing, the tree-level Stasheff
 at arity $2$);
\item $\h^1$, one loop:
 $d_{\mathrm{sew}}^2 + [d_{\mathrm{int}},\,
 \h\Delta_{\mathrm{BV}}] = 0$ (the one-loop QME, producing
 the curvature $\kappa$);
\item $\h^1$, two edges:
 $d_{\mathrm{sew}}\,\h\Delta_{\mathrm{BV}}
 + \h\Delta_{\mathrm{BV}}\,d_{\mathrm{sew}} = 0$
 (compatibility of sewing with the BV Laplacian);
\item $\h^2$: $(\h\Delta_{\mathrm{BV}})^2 = 0$ (the BV
 Laplacian squares to zero, reflecting the second-order
 nature of $\Delta_{\mathrm{BV}}$).
\end{itemize}
Each relation is a Stokes identity on the corresponding stratum of
the modular operad.
\end{proof}

\subsection{The universal MC element and the genus expansion}
% label removed: subsec:thqg-bv-ext-mc-genus

\begin{theorem}[BV amplitudes $=$ MC projections;
\ClaimStatusProvedHere]
% label removed: thm:thqg-bv-ext-mc-projections
The total amplitude
$\cA = \sum_{g,k} \h^g\,\cA_k^{(g)}$
is the image under the stable-graph expansion of the universal
Maurer--Cartan element
\[
\Theta_A = D_A - d_0
\;\in\; \mathrm{MC}\bigl(\Def_{\mathrm{cyc}}^{\mathrm{mod}}(A)
\;\hat\otimes\; \cG_{\mathrm{mod}}\bigr),
\]
where $\Def_{\mathrm{cyc}}^{\mathrm{mod}}(A)$ is the modular
cyclic deformation complex and $\cG_{\mathrm{mod}}$ is the
stable-graph coefficient algebra (Vol~I,
Definition~\ref{V1-def:stable-graph-coefficient-algebra}).
Specifically:
\begin{enumerate}[label=(\roman*)]
\item At genus $0$, projection to arity $k$ gives
 $m_k = \Theta_A^{(0,k)}$: the tree-level $\Ainf$ operations;
\item At genus $0$, arity $2$: $\kappa = \Theta_A^{(0,2)}$ is
 the curvature (the binary scalar shadow);
\item At genus $1$, $\Theta_A^{(1,k)}$ gives the one-loop
 amplitude $\cA_k^{(1)}$ on $\overline{\cM}_{1,k}$;
\item At genus $g$, $\Theta_A^{(g,k)} = \cA_k^{(g)}$ on
 $\overline{\cM}_{g,k}$.
\end{enumerate}
The Maurer--Cartan equation
$D\,\Theta_A + \tfrac{1}{2}[\Theta_A, \Theta_A] = 0$
encodes simultaneously: the $\Ainf$ Stasheff identities (genus $0$),
the curved bar differential (genus $1$), the three-shell identity
(genus $2$), and all higher-genus Stokes identities.
\end{theorem}

\begin{proof}
The identification of $\Theta_A = D_A - d_0$ as a Maurer--Cartan
element is proved in Vol~I
(Theorem~\ref*{V1-thm:mc2-bar-intrinsic}):
$\Theta_A$ is MC because $D_A^2 = 0$ (the square of the total
differential on the bar complex vanishes, being the
Feynman-transform differential of a modular operad algebra).
The $(g,k)$-projection of the MC equation is exactly the
Stokes identity on $\overline{\cM}_{g,k}$ applied to the
weight forms of the BV theory:
\[
\bigl(D\,\Theta_A + \tfrac{1}{2}[\Theta_A, \Theta_A]
\bigr)^{(g,k)} = 0
\quad\Longleftrightarrow\quad
\int_{\overline{\cM}_{g,k}} d\,\omega^{(g)}_k = 0.
\]
The left-hand side decomposes into the three components
$d_{\mathrm{int}}$, $d_{\mathrm{sew}}$, and
$\h\Delta_{\mathrm{BV}}$ of
Corollary~\ref{cor:thqg-bv-ext-three-components}, each
producing the corresponding boundary stratum of
$\overline{\cM}_{g,k}$. The bracket $[\Theta_A, \Theta_A]$
contributes the sewing and BV-Laplacian terms (which are
quadratic in $\Theta_A$ via graph-gluing), while the linear
$D\,\Theta_A$ contributes the internal differential.

The identification at genus $0$ with the tree-level $m_k$ is
Theorem~\ref{thm:thqg-bv-ext-trees-bar}.
The identification at genus $1$ with the curvature $\kappa$ is
Theorem~\ref{thm:thqg-bv-ext-kappa}.
The identification at genus $2$ with the three-shell identity is
Theorem~\ref{thm:thqg-bv-ext-three-shell}.
The passage to arbitrary genus is by the structural identity of
the Feynman transform (Getzler--Kapranov): the graph sum defining
$\Theta_A$ is precisely $\mathsf{F}(\mathsf{Com})$ applied to $A$.
\end{proof}

\begin{remark}[The Chriss--Ginzburg principle]
% label removed: rem:thqg-bv-ext-cg-principle
Theorem~\ref{thm:thqg-bv-ext-mc-projections} is an instance of
the organizing principle of this monograph: every algebraic
structure is a Maurer--Cartan element in a convolution dg Lie
algebra. The $\Ainf$ structure ($m_k$), the curvature ($\kappa$),
the three-shell identity, and the full genus expansion are all
projections of the single element $\Theta_A$. The BV-BRST
derivation of this chapter constructs $\Theta_A$ from the physics
side (Feynman diagrams on $\FM_k(\bC) \times \Conf_k(\bR)$ and
stable-graph amplitudes on $\overline{\cM}_{g,k}$), while Vol~I
constructs it from the algebra side (the bar-intrinsic definition
$\Theta_A = D_A - d_0$). The two constructions agree by the
identification of Theorem~\ref{thm:thqg-bv-ext-modular-bar}.
\end{remark}

\begin{remark}[Shadow projections from BV data]
% label removed: rem:thqg-bv-ext-shadow-projections
The shadow obstruction tower of Vol~I (the finite-order truncations
$\Theta_A^{\le r}$) admits a direct BV interpretation.
At arity $r = 2$: $\kappa = \operatorname{str}(L_0)$ is the
one-loop vacuum tadpole. At arity $r = 3$: the cubic shadow
$C_A$ is the three-point one-loop amplitude on
$\overline{\cM}_{1,3}$. At arity $r = 4$: the quartic resonance
class $Q_A$ is the four-point amplitude on $\overline{\cM}_{1,4}$,
with the clutching law arising from the boundary of
$\overline{\cM}_{1,4}$ (node formation on the torus). The
shadow depth classification (G/L/C/M) is the classification of
theories by the arity at which their BV amplitudes first receive
nonvanishing contributions from nonlinear vertices: Gaussian
theories ($r_{\max} = 2$) have only tadpoles, Lie-type theories
($r_{\max} = 3$) have nonvanishing three-point amplitudes, and
so on.
\end{remark}

\subsection{Convergence and the sewing envelope}
% label removed: subsec:thqg-bv-ext-convergence

\begin{proposition}[HS-sewing from BV data; \ClaimStatusProvedHere]
% label removed: prop:thqg-bv-ext-hs-sewing
Under the Hilbert--Schmidt sewing condition
\begin{equation}
% label removed: eq:thqg-bv-ext-hs-condition
\sum_{a,b,c \ge 0} q^{a+b+c}\,
\|m^c_{a,b}\|_{\mathrm{HS}}^2 < \infty,
\end{equation}
where $m^c_{a,b} : A_a \otimes A_b \to A_c$ are the
weight-graded components of $m_2$ and $\|-\|_{\mathrm{HS}}$ is
the Hilbert--Schmidt norm, the stable-graph amplitudes
$\cA_k^{(g)}$ converge absolutely for all $g$ and $k$, and the
genus expansion
\[
\cA_k = \sum_{g \ge 0} \h^g\,\cA_k^{(g)}
\]
defines a well-defined element of $A^{\otimes k}\llbracket\h\rrbracket$.
\end{proposition}

\begin{proof}
The HS condition~\eqref{eq:thqg-bv-ext-hs-condition} ensures that
the sewing operator $\Delta_{\mathrm{sew}} = \sum_n a_n \otimes
a^n$ converges in the trace-class sense: the contraction at each
edge $e$ of a stable graph $\Gamma$ involves
$\Tr(\Delta_e) = \sum_n \langle a_n, a^n \rangle$, which converges
by the HS condition. The product over all edges of $\Gamma$ is
bounded by
\[
|\cA_\Gamma| \le \prod_{e \in E(\Gamma)}
\|\Delta_e\|_{\mathrm{tr}}
\cdot \prod_{v \in V(\Gamma)} C_{g_v, n_v},
\]
where $C_{g_v, n_v}$ is a constant depending on the geometry of
$\overline{\cM}_{g_v, n_v}$ (bounded by the volume of the moduli
space times the sup-norm of the vertex weight form).

For the standard landscape (polynomial OPE growth and
subexponential sector growth), the HS condition is verified in
Vol~I (Theorem~\ref{thm:general-hs-sewing}). The
graph sum over $\mathrm{Graph}(g,k)$ converges because the number
of stable graphs of type $(g,k)$ grows at most exponentially in
$g$, while the HS bounds provide exponential decay from the sewing
operators. The detailed estimates are in
Section~\ref{sec:thqg-bv-ext-one-loop} for the one-loop case and
extend to all genera by induction on the boundary structure of
$\overline{\cM}_{g,k}$.
\end{proof}

\begin{remark}[Algebraic vs.\ analytic bar]
% label removed: rem:thqg-bv-ext-algebraic-analytic
The algebraic bar complex $B(A)$ with its modular differential
$D_{\mathrm{mod}}$ is an algebraic object: it is defined for any
chiral algebra $A$ without convergence hypotheses. The analytic
bar complex $B^{\mathrm{an}}(A)$ (the graph-norm closure of $B(A)$
under the HS topology) is the analytic completion required for
sewing and partition functions. The HS condition bridges the two:
it ensures that the algebraic structure maps extend continuously to
the analytic completion. The sewing envelope $A^{\mathrm{sew}}$
(the Hausdorff completion for all-amplitude seminorms) sits between
the algebraic $A$ and the Hilbert-space completion
$\overline{A}$: we have $A \subset A^{\mathrm{sew}} \subset
\overline{A}$, with $B^{\mathrm{an}}(A)$ computing the correct
genus-$g$ amplitudes for all $g \ge 0$.
\end{remark}

%% ====================================================================
%% CODA
%% ====================================================================

\begin{remark}[Summary of the BV-to-bar dictionary]
% label removed: rem:thqg-bv-ext-dictionary
The results of this chapter establish the following dictionary
between BV-BRST data and bar-complex algebra:

\medskip
\begin{center}
\renewcommand{\arraystretch}{1.3}
\begin{tabular}{l|l}
\textbf{BV-BRST (physics)} & \textbf{Bar complex (algebra)} \\
\hline
QME: $\h\Delta_{\mathrm{BV}}S + \frac{1}{2}\{S,S\} = 0$
 & $D_{\mathrm{mod}}^2 = 0$ \\
BRST operator $Q = \{S,-\}$
 & $d_{\mathrm{int}} = m_1$ \\
Propagator $G(z_1,z_2;t_1,t_2)$
 & Logarithmic $1$-form $\eta_{ij}$ \\
Tree-level Feynman diagrams
 & $m_k = \sum_{\mathrm{Tree}(k)} |\Aut|^{-1}\omega_\Gamma$ \\
BV Laplacian $\Delta_{\mathrm{BV}}$
 & Self-loop creation $\h\Delta_{\mathrm{BV}}$ \\
One-loop tadpole
 & Curvature $\kappa = \operatorname{str}(L_0)$ \\
$g$-loop Feynman diagrams
 & Stable-graph amplitude $\cA_k^{(g)}$ \\
Boundary of $\overline{\cM}_{g,k}$
 & Stasheff identity at genus $g$ \\
$\partial^2 = 0$ on $\overline{\cM}_{g,k}$
 & Corner cancellation (AOS) \\
Full QME genus expansion
 & MC element $\Theta_A \in
 \mathrm{MC}(g^{\mathrm{mod}}_A)$ \\
HS-sewing condition
 & Analytic bar convergence
\end{tabular}
\end{center}

\medskip\noindent
The $\bC$-direction of the FM product space governs the bar
differential (holomorphic factorization); the $\bR$-direction
governs the bar coproduct (topological factorization). Together
they form the Swiss-cheese algebra
$\SCchtop$ of Section~\ref{subsec:from-prefact},
with the BV-BRST derivation of this chapter providing its
perturbative construction.
\end{remark}
